\section{讨论与限制}

下面五条加本篇各一条，构成本篇的全部已知边界。

\begin{itemize}
\item \textbf{L1 同源自证，这是首要威胁。} 关系表与人物表出自同一部原书，被统计的叙事与关系标记可能是同一段文字的两种编码。处理：把它写进摘要而非脚注，且不把 $\Delta_c$ 叫效应量，只叫相对基线的位移。
\item \textbf{L2 短词过度匹配。} 1--2 字主题词逐字命中会把同字根的无关词一并按住，排序对匹配规则敏感。处理：把这类高频词降到定位词，论据只放在有对称性的比较上。
\item \textbf{L3 类别是粗标签。} 关系类别里通常有一类「其他」占四成上下且几乎不动，没有区分力，只当分母；真正有区分力的是少数几类，且它们方向相反。类别边界漂移最可能落在这里。本篇子图「其他」类占 49.1，基线 39.9；聚焦并没有换来更好的编码分辨率。
\item \textbf{L4 子集小，一个小数就是一大步。} 本子图显式边 636 条，单类位移常由两三条边贡献；一条被重标类别，比例就跳一两成。处理：只报符号与量级，并在正文表里把边数一并写出。
\item \textbf{L5 自审不是跨模型对抗评审。} 本机无第二模型端点，评审由同一模型多轮盲读完成。已知后果是写作轮曾把未核验的文献键写进正文，盲读轮才拦下。自审能抓内部不一致，抓不到共同先验盲区。
\item \textbf{L6 锚点是人工选定的，而且只有四个。} 口径 B 的四个锚点取自该母题的经典人物
（朱元璋、朱棣、朱由检、方孝孺），不是随机抽样；换一套锚点，$\Delta_c$ 的符号未必保持，
本篇因此也不声称这个子图能代表「明代宗族」一类的总体。
\item \textbf{L7 母题只在词汇与人物关系两层被验证。} 本篇主题事件仅 3 件（1404、1581、1601 各 1 件）、
地点 0 处，「家国同构」没有落到可定位的事件序列或空间分布上做第三层检验；
任何关于这一母题随时间演变的断言，都超出本篇证据。
\end{itemize}

需要强调的是：以上限制没有一条是可以通过再跑一遍数据消掉的。它们不是工程瑕疵，而是这类「从单一叙事抽取结构化图谱再回看该叙事」的研究在方法上的固有位置。本文的处理是把它摊开：能降到位移的，就不叫效应；不能降的，就明确标为自证循环。
